\documentclass[a4paper,12pt]{article}
\usepackage[utf8]{inputenc}
\usepackage{cmap}
\usepackage[T2A]{fontenc}
\usepackage[english,russian]{babel}
\usepackage{amsmath}
\usepackage{amsfonts}
\usepackage{amssymb}
\usepackage{geometry}
\geometry{a4paper, top=2cm, bottom=2cm, left=2.5cm, right=2.5cm}
\usepackage{enumitem}
\setlength{\parindent}{0pt}
\setlength{\parskip}{0.7em}

\begin{document}

\begin{center}
\Large\textbf{Машинное обучение 1}
\end{center}
\begin{center}
\large Контрольная работа (2021 год) --- Вариант 1
\end{center}
\hrule
\vspace{1cm}

\subsection*{Задача 1 (2.5 балла)}

\textbf{1.} Метод one-vs-all (один против всех) для многоклассовой классификации:
Для каждого класса \(k = 1,\dots,K\) обучается бинарный классификатор \(a_k(x)\), который предсказывает \(+1\), если объект относится к классу \(k\), и \(-1\) иначе.
При классификации нового объекта выбирается класс, давший максимальный ответ:
\[
\hat{y}(x) = \arg\max_{k=1,\dots,K} a_k(x)
\]

\bigskip

\textbf{2.} Рассмотрим функции потерь для \(y = -1\):
\begin{itemize}
  \item \(L_a(-1, z) = \log(1 + \exp(z))\). Функция монотонно убывает при \(z \to -\infty\), минимум достигается асимптотически при \(z \to -\infty\), поэтому минимальное по модулю значение \(z\) не существует (или равно \(+\infty\))
  \item \(L_b(-1, z) = \max(0, 1 + z)\). Минимум достигается при \(z \le -1\), минимальное по модулю \(z = -1\), \(|z| = 1\)
  \item \(L_c(-1, z) = (-z - 1)^2 = (z + 1)^2\). Минимум при \(z = -1\), \(|z| = 1\)
\end{itemize}

\bigskip

\textbf{3.} Добавим порог \(t = 10\): \(\hat{a}(x) = \operatorname{sign}(\langle w, x \rangle - 10)\).
Для положительных объектов (\(y = +1\)) исходный классификатор даёт \(\langle w, x \rangle > 0\).
Добавление порога может изменить знак, если \(0 < \langle w, x \rangle < 10\).
\begin{itemize}
  \item Для \(L_a\) значение может измениться, так как функция потерь зависит от отступа \(y \cdot \langle w, x \rangle\)
  \item Для \(L_b\) при больших отступах (\(\langle w, x \rangle \ge 1\)) значение равно 0, и порог его не изменит, если отступ останется \(\ge 1\), то есть \(\langle w, x \rangle \ge 11\). В общем случае утверждение неверно
\end{itemize}

\bigskip

\textbf{4.} SVM с \(C = 1\) относит все объекты к синему классу, вероятно, из-за сильного дисбаланса классов: синий класс составляет большинство, и штраф за ошибку на красных объектах перекрывается выигрышем от правильной классификации синих. Увеличение параметра \(C\) может помочь сбалансировать штрафы за ошибки на разных классах

\newpage

\subsection*{Задача 2 (2.5 балла)}

\textbf{1.} Out-of-bag (OOB) оценка в бэггинге:
При построении каждого дерева с помощью бутстрапа часть объектов (в среднем около 37\%) не попадает в выборку и является "контрольной" для данного дерева.
OOB-ошибка вычисляется как средняя ошибка предсказаний по тем деревьям, которые не видели объект при обучении.
Используется для оценки качества модели без отдельной отложенной выборки и для настройки гиперпараметров

\bigskip

\textbf{2.} Ошибки в формулах:
\begin{itemize}
  \item В определении \(s_i\) должна быть производная по \(a_{N-1}(x_i)\), а не по \(b_{N-1}(x_i)\)
  \item Должен быть знак минус: \(s_i = -\frac{\partial L(y_i, z)}{\partial z}\big|_{z = a_{N-1}(x_i)}\)
  \item В целевой функции для \(b_N\) обычно используют среднеквадратичную ошибку: \(\sum (b_N(x_i) - s_i)^2\), а не \(\frac1\ell \sum L(s_i, b_N(x_i))\)
\end{itemize}

\bigskip

\textbf{3.} Для регрессии с MSE:
\begin{itemize}
  \item Решающее дерево: прогноз в листе — среднее целевых значений. Если все \(y_i > 0\), то среднее \(> 0\), поэтому прогнозы неотрицательны
  \item Случайный лес: то же самое, прогнозы неотрицательны
  \item Градиентный бустинг: базовые деревья обучаются на сдвигах (псевдо-остатках), которые могут быть отрицательными, поэтому прогнозы отдельных деревьев и итоговой композиции могут быть отрицательными
\end{itemize}

\bigskip

\textbf{4.} В XGBoost регуляризация нормы весов в листьях (\(\frac{\lambda}{2}\sum b_j^2\)) работает, потому что веса подбираются аналитически при фиксированной структуре дерева.
В одиночном дереве, обучаемом обычным жадным алгоритмом, веса в листьях определяются как среднее целевых значений, и регуляризация нормы не вписана в процесс обучения.
Её добавление потребовало бы пересмотра критерия разбиения и метода подбора прогнозов

\newpage

\subsection*{Задача 3 (2.5 балла)}

\textbf{1.} Логистическая функция потерь гладкая и выпуклая, её градиент хорошо определён везде, что позволяет использовать градиентные методы оптимизации.
Доля ошибок кусочно-постоянна, её градиент равен нулю почти везде, что делает оптимизацию практически невозможной с помощью градиентного спуска

\bigskip

\textbf{2.}
\begin{itemize}
  \item \textit{Может ли доля ошибок вырасти, а логистическая ошибка упасть?}
        Да, это возможно: модель может стать увереннее в правильных ответах, но один объект перестал классифицироваться правильно
  \item \textit{Может ли логистическая ошибка вырасти, а доля ошибок упасть?}
        Тоже возможно: например, модель уменьшила уверенность в правильных ответах, но при этом все ответы остались верными
\end{itemize}

\bigskip

\textbf{3.} Пример функции потерь:
\[
L(y,a) = 
\begin{cases}
\frac12 (y-a)^2, & a \le y,\\
|y-a|, & a > y
\end{cases}
\]
Штрафует за завышение прогноза квадратично, а за занижение — линейно, что обеспечивает устойчивость к выбросам при занижении и сильный штраф при завышении.
Однако она не везде дифференцируема.
Гладкий вариант: \(L(y,a) = \frac12 (y-a)^2 + \frac12 \log\cosh(a-y)\) с возможной настройкой коэффициентов для асимметрии

\bigskip

\textbf{4.}
\begin{itemize}
  \item \(L_3 = \sum |w_j|^3\): штрафует большие веса сильнее, чем \(L_2\), и стремится сделать веса равномерно маленькими. Может давать более разреженные решения, чем \(L_2\), но менее разреженные, чем \(L_1\)
  \item \(L_\infty = \max_j |w_j|\): стремится минимизировать максимальный по модулю вес, что приводит к выравниванию весов. Не дифференцируема в точках, где два веса равны по модулю
  \item \(L_0 = \sum [w_j \neq 0]\): число ненулевых весов. Даёт максимально разреженное решение, но задача оптимизации становится NP-трудной (комбинаторная)
\end{itemize}

\newpage

\subsection*{Задача 4 (2.5 балла)}

Дано:
\[
x_i \sim \mathcal{N}(0, \sigma^2), \quad y_i = \varepsilon_i x_i, \quad \varepsilon_i \sim \text{Exp}(\lambda),
\]
все \(\varepsilon_i\) и \(x_i\) независимы.
Известно: \(\mathbb{E}[\varepsilon_i] = \frac{1}{\lambda},\quad \mathbb{D}[\varepsilon_i] = \frac{1}{\lambda^2}\)

Модель:
\[
\mu(X)(x) = \frac{1}{\lambda} \left( \sum_{i=1}^{\ell} x_i \right) \operatorname{sign}(x)
\]

\bigskip

\textbf{1. Найдём \(\mathbb{E}[y \mid x]\):}
\[
\mathbb{E}[y \mid x] = \mathbb{E}[\varepsilon x \mid x] = x \cdot \mathbb{E}[\varepsilon] = \frac{x}{\lambda}
\]

\bigskip

\textbf{2. Шум (noise):}
\[
\text{Noise} = \mathbb{E}_{x,y}\bigl[(y - \mathbb{E}[y \mid x])^2\bigr] = \mathbb{E}_x\bigl[x^2 \cdot \mathbb{D}[\varepsilon]\bigr] = \mathbb{E}_x[x^2] \cdot \frac{1}{\lambda^2} = \sigma^2 \cdot \frac{1}{\lambda^2}
\]

\bigskip

\textbf{3. Средний ответ модели:}
\[
\mathbb{E}_X[\mu(X)(x)] = \frac{1}{\lambda} \cdot \mathbb{E}_X\Bigl[\sum_{i=1}^{\ell} x_i\Bigr] \cdot \operatorname{sign}(x) = \frac{1}{\lambda} \cdot 0 \cdot \operatorname{sign}(x) = 0
\]

\bigskip

\textbf{4. Смещение (bias):}
\[
\text{Bias}^2 = \mathbb{E}_x\bigl[(\mathbb{E}_X[\mu(X)(x)] - \mathbb{E}[y \mid x])^2\bigr] = \mathbb{E}_x\bigl[(0 - x/\lambda)^2\bigr] = \frac{1}{\lambda^2} \mathbb{E}_x[x^2] = \frac{\sigma^2}{\lambda^2}
\]

\bigskip

\textbf{5. Разброс (variance):}
\[
\text{Var} = \mathbb{E}_x\bigl[\mathbb{D}_X(\mu(X)(x))\bigr] = \mathbb{E}_x\Bigl[\operatorname{sign}^2(x) \cdot \frac{1}{\lambda^2} \cdot \mathbb{D}_X\Bigl(\sum_{i=1}^{\ell} x_i\Bigr)\Bigr] = \frac{1}{\lambda^2} \cdot 1 \cdot \ell \sigma^2 = \frac{\ell \sigma^2}{\lambda^2}
\]

\bigskip

Итоговые компоненты разложения ошибки:
\[
\boxed{\text{Noise} = \frac{\sigma^2}{\lambda^2}, \quad \text{Bias}^2 = \frac{\sigma^2}{\lambda^2}, \quad \text{Variance} = \frac{\ell \sigma^2}{\lambda^2}}
\]

\end{document}
